% Lösungen Rekonstruktion von Beständen (zusammenbau v0.4, Heft)
\einheitenkopf{Rekonstruktion von Beständen}

\erg{22}{a) eine Rate – Liter je Minute ist eine Menge pro Zeit \quad b) $[t^3 - 6t^2 + 15t]_0^4 = 28$; in den ersten vier Stunden fließen 28 Liter zu \quad c) $50 + ∫_0^3 r(t)\,\mathrm{d}t = 50 + [2t^2 + 2t]_0^3 = 50 + 24 = 74$ Liter}
\erg{23}{a) Dreieck unter $c'$: $\tfrac{1}{2} \cdot 8 \cdot 5 = 20$; $c(0) = 90 - 20 = 70$ mg/l (die Fläche wird abgezogen) \quad b) Änderung $\approx 4$\,°C (Fläche über der Achse von 0 bis 4); die Rate ist dort positiv, die Temperatur nimmt zu}
\erg{24}{a) $t = 5$ – dort wechselt r von plus nach minus \quad b) Nullstellen von r: $t = 0$ und $t = 8$; dazwischen ist $r(t) > 0$ (etwa $r(4) = 8$), also nimmt der Inhalt durchgehend zu \quad c) $t = 4$: dort ist $r(t) = 0$, und r wechselt von plus nach minus; vorher wächst, danach fällt der Inhalt}
\erg{25}{a) Nullstellen von r: $t = 0$ und $t = 6$; dazwischen $r(t) > 0$, also nimmt das Volumen zu – ob r dort steigt oder fällt, spielt keine Rolle \quad b) falsch: bei $t = 8$ hat r den Hochpunkt, dort wächst der Inhalt am schnellsten; $r(t) = t^2 \cdot (3 - 0{,}25t)$ wird erst bei $t = 12$ null und wechselt dort von plus nach minus – am meisten ist bei $t = 12$ im Silo}
\erg{26}{a) $r(1) = 3$: in diesem Moment fließen 3 Liter je Minute zu; Dreiecksfläche von 0 bis 4: Änderung $= 8$ Liter (Zunahme) \quad b) $∫_0^{10} (20 - t)\,\mathrm{d}t = 200 - 50 = 150$; nach zehn Sekunden liegt das erste Auto 150\,m vorn}
\erg{27}{a) $∫_0^4 (6t - 1{,}5t^2)\,\mathrm{d}t = [3t^2 - 0{,}5t^3]_0^4 = 48 - 32 = 16$; nach vier Wochen ist die erste Pflanze 16\,cm höher als die zweite (nicht 16\,cm hoch) \quad b) z ist der Zeitpunkt, zu dem beide gleich weit gelaufen sind (Streckendifferenz null). Bis zum Schnittpunkt wächst der Vorsprung der ersten Läuferin, erst danach wird er abgebaut; null kann er also erst nach dem Schnittpunkt wieder werden. \quad c) $6z - 0{,}25z^2 = 0$, also $z = 24$: nach 24 Sekunden sind beide gleich weit gefahren. Die Geschwindigkeiten sind bei $t_s = 12$ gleich; $z$ liegt danach, weil der Vorsprung bis $t_s$ wächst und erst danach abgebaut wird}
